% from ega4-2-fr.tex
\oldpage[IV-2]{31}
il en est de même de $g^*(\theta)$, et comme $g$ est fidèlement plat,
$\theta$ est bijectif (2.2.7), ce qui prouve (viii).

\emph{(ix)} La proposition résulte de (viii), de (I, 5.3.4) et de (I, 5.3.8)
qui ramène les monomorphismes aux isomorphismes.

\emph{(x)} Si $f'$ est une immersion ouverte, $f'(X')$ est ouvert dans $Y'$,
et l'on a $f'(X')=g^{-1}(f(X))$ (I, 3.4.8)~; il résulte de (2.3.12) que
$f(X)$ est ouvert. On peut alors remplacer $Y$ (resp. $Y'$) par le
sous-préschéma induit sur l'ouvert $f(X)$ (resp. $f'(X')$), compte tenu de
(1.1.2) et (2.2.13)~; alors $f'$ devient un isomorphisme, donc il en est de
même de $f$ par (viii), et cela établit (x).

\emph{(xi)} Si $f'$ est une immersion quasi-compacte, $f'$ est un morphisme
quasi-compact et quasi-séparé (1.2.2), donc il en est de même de $f$ par (ii)
et (2.6.2, (v)). Soit $Z$ le sous-préschéma de $Y$ image fermée de $X$ par
$f$ (1.7.8), et posons $f=j\circ q$, où $j:Z\to Y$ est l'injection
canonique~; on a alors $f'=j'\circ q'$ avec $j'=j_{(Y')}$,
$q'=q_{(Y')}$, et l'on sait que $j'$ s'identifie à l'injection canonique
$Z'\to Y'$ du sous-préschéma $Z'$ de $Y'$, image fermée de $X'$ par $f'$
(2.3.2). L'hypothèse sur $f'$ signifie alors que $q'$ est une immersion
ouverte (I, 9.5.10), donc il en est de même de $q$ par (x), et cela montre
que $f$ est une immersion.

\emph{(xii)} Dire que $f$ (resp. $f'$) est une immersion fermée signifie que
$f$ (resp. $f'$) est une immersion quasi-compacte et un morphisme fermé~; on
voit donc que (xii) résulte de (xi) et de (2.6.2, (ii)).

\emph{(xiii)} et \emph{(xiv)} Supposons $f'$ affine (resp. quasi-affine)~;
notons alors que $f'$ est quasi-compact et quasi-séparé (II, 5.1.1), donc il
en est de même de $f$ par (ii) et (2.6.2, (v)). Posons
$\mathcal A=f_*(\mathcal O_X)$, $\mathcal A'=f'_*(\mathcal O_{X'})$~; en
vertu de (2.3.1), l'homomorphisme canonique de $\mathcal O_{Y'}$-Algèbres
$g^*(\mathcal A)\to\mathcal A'$ est bijectif~; par conséquent, si
$h:Z=\operatorname{Spec}(\mathcal A)\to Y$ est le morphisme structural, le
morphisme structural $h':Z'=\operatorname{Spec}(\mathcal A')\to Y'$
s'identifie à $h_{(Y')}$ (II, 1.5.2). Soit alors $u:X\to Z$ (resp.
$u':X'\to Z'$) le $Y$-morphisme (resp. $Y'$-morphisme) canonique
correspondant à l'homomorphisme identique de $\mathcal A$ (resp.
$\mathcal A'$) (II, 1.2.7)~; comme on a le diagramme commutatif
\[
\xymatrix{
X \ar[d]_{u} & X' \ar[l] \ar[d]^{u'}\\
Z \ar[d]_{h} & Z' \ar[l]_{g'} \ar[d]^{h'}\\
Y & Y' \ar[l]^{g}
}
\]
et que $h'\circ u'=f'$, il résulte de (II, 1.2.7) que $u'=u_{(Z')}$. En
outre, $g'$ est fidèlement plat et quasi-compact (1.1.2 et 2.2.13). Cela
étant, l'hypothèse sur $f'$ signifie que $u'$ est un isomorphisme (resp. une
immersion ouverte) (II, 5.1.6)~; il résulte alors de (viii) (resp. (x)) que
$u$ est un isomorphisme (resp. une immersion ouverte), d'où (xiii) (resp.
(xiv)).

